% Generated by scripts/export_feature_tables.py -- do not edit by hand.
\begin{description}
  \item[\texttt{item\_popularity\_count}] How often an item was viewed in training sessions up to the cut-off. Drives the popularity baseline and the cold-start fallback ordering for items with no neighbour-session support.

  \item[\texttt{item\_category}] Catalogue attribute from the item-properties changelog. Known before any interaction, so it is available for items that never appear in training and can be used to break ties among equally popular candidates.

  \item[\texttt{visitor\_prior\_session\_count}] Number of training sessions a visitorid has started before the current one. This is an optional, strictly ablation-only signal: visitorid is a cookie, not a durable account, so the primary session-kNN method never depends on it, but a returning cookie is still weak positive evidence.

  \item[\texttt{session\_length\_so\_far}] Number of events already observed in the current session, i.e. the 0-indexed position itself. Short sessions get little context, which is exactly the "cold session" regime broken out in the evaluation report.

  \item[\texttt{time\_since\_session\_start}] Elapsed time between the session's first event and the *last observed* event (not the event being predicted). Captures browsing pace -- a visitor who has lingered for ten minutes behaves differently from one three seconds into a session.

  \item[\texttt{hour\_of\_day}] Hour extracted from the event being scored. Not an aggregate over history, so it carries no leakage risk: the wall-clock time of the prediction request is known by definition.

  \item[\texttt{is\_weekend}] Weekday/weekend split of browsing behaviour, used as an ablation feature alongside hour\_of\_day to test whether coarse temporal context improves candidate re-ranking over neighbour-vote scores alone.

\end{description}
